\chapter{Relative Velocity and the Sign of the Interaction}

The previous chapter turned the signal off. It stated a stop rule and obeyed it: the
account owed a reading to a requested depth, it paid that reading, and it declined to
add another layer merely because more could be said. That rule still holds, and this
chapter does not break it. What follows is not a deeper cut at the same place. It is a
different measurement of something the produced pair has carried since the moment it
was produced, a reading the apparatus had not yet taken because nothing had yet asked
for it. The stop forbids more depth. It does not forbid reading what was already there.

What was already there is a relative velocity. The single invariant became a charge, and
the charge came as a signed pair; the apparatus has spoken of that pair as an orientation
of a residue, read once forward and once backward. Read instead as a motion, the same
residue is the rate at which the two parties of the pair separate from one another --- a
relative velocity. The number does not change. The orientation that was a sign becomes the
sign of a velocity, and the question the book postponed about that sign comes back in a
sharper form: the sign of \emph{whose} velocity? The answer is the chapter's first claim,
and it is not the obvious one. The relative velocity is not a property of either party. It
is a property of the interaction between them, and only the interaction carries it.

The first section re-reads the pair as a relative velocity and shows that the reading has
the same shape every measurement in this book has had: a whole against its parts, with the
remainder as the residue. The coupled reading of the moving pair is the whole; the two
one-sided stories, each party told on its own, are the parts; and the relative velocity is
what is left when the parts are subtracted from the whole. That remainder is the single
invariant under a new name. It is read, as everything here is read, and not constructed.

The second section makes precise why the velocity belongs to the interaction. A
one-sided story knows nothing of relative motion; relative velocity is a cross term, a
reading of the two parties against each other, and a single coordinate is blind to it
exactly as the interior energy was blind to the sign. The magnitude of the motion is real
and finite, but its sign lives where the interior cannot reach: in the term that requires
both parties at once. To read the sign, the apparatus must hold the pair, not either half.

The third section confronts a quieter problem and answers it with the most disciplined move
in the book. The apparatus has two ways to reach the relative velocity --- it can assemble
the reading from measurements of the path, or it can read the invariant directly --- and the
two ways return the same number. That they return the same number is a finite fact and can
be checked. But for the apparatus to \emph{treat} the two readings as one reading, it must
issue a single receipt to that effect: one act of identification, located and named, by
which two records of the same velocity are collapsed into a single thing. The book has made
countless routine identifications; this is not one of them. It is the one genuine selection
the construction performs, and the third section isolates it so the reader can see it
standing alone.

The fourth section returns to the sign and to the convention. The relative velocity carries
a sign; the sign splits the pair into two orientations that cancel; and which orientation is
named \emph{matter} is not read off the velocity at all. It is fixed by counting --- the same
tally that anchored the convention before. The count is the measurement; the name is the
convention laid on top of it. The fifth section states the ceiling and refuses to cross it.
The split is the sign of the relative-velocity invariant, measured by the apparatus; the
orientation is a convention; and that is the whole of the claim. No asymmetry is derived,
no imbalance of the early universe is explained, no mechanism is proposed. The pair, summed,
still cancels, and the cancellation is structural, not a result.

It helps to carry one figure through the chapter. Picture the apparatus as a nest of
enclosures. The outermost holds the measurement; the measurement holds the coupling of the
two parties; the coupling holds the two parties, one within the other; and at the center,
where one more enclosure would close the figure, the apparatus draws none. The innermost
bracket is left open:
\[
  \underbrace{\Big(\,}_{\text{apparatus}}
  \underbrace{\Big(\,}_{\text{measurement}}
  \underbrace{\big(\,}_{\text{coupling}}
  \,x\,\big(\,y\,\big(\,\cdots
\]
The open bracket is not a slip of the pen and not an omission to be repaired. It is the
residue. It is the one quantity the enclosures reach but cannot close --- the difference
between the coupled measurement and the stories that would explain it away --- and the
apparatus reads it precisely there, at the place the construction leaves un-closed. Every
section below returns to that open bracket. The relative velocity lives in it; the sign is
the orientation in which it is read; and the honesty of the chapter is the honesty of an
instrument that reads an open quantity and never pretends to have closed it.

\section{The Pair Re-Read as Relative Velocity}

The pair was produced as a residue. Two coordinates were kicked together, the apparatus took
a coupled reading of the result, and it compared that reading against the linear story --- the
account in which each coordinate moves on its own and nothing couples. The residue was the
difference: what the coupled measurement showed that the one-sided stories could not. The
earlier chapters read that residue as a charge and as an orientation. This section reads it as
a motion, and finds the same number.

A relative velocity is a reading of how two things move with respect to each other. It is not
the motion of the first thing, recorded against a fixed background, nor the motion of the
second; it is the difference of the two motions, the rate at which the gap between them opens
or closes. The pair offers exactly such a reading. The coupled measurement records the two
parties as they actually moved, together; the two stories record each party as it would have
moved alone. Subtract the stories from the coupled measurement and what remains is the part of
the motion that only the pair has --- the relative velocity.

The arithmetic is the book's oldest arithmetic. Call the coupled reading the whole and the two
one-sided variations the parts. The relative velocity is the whole minus the parts:
\[
  \text{relative velocity} \;=\; \text{whole} \;-\; \text{parts}.
\]
This is not a new quantity smuggled in under a new word. The whole, expanded, is exactly the
two one-sided variations plus a remainder, and the remainder is the single invariant the book
has measured since it first forced a unique second-order term. So the subtraction leaves the
invariant standing alone:
\[
  \text{relative velocity} \;=\; \big(\text{parts} + \text{invariant}\big) - \text{parts}
  \;=\; \text{invariant}.
\]
The pair's relative velocity is the single invariant. The charge, the orientation, and the
relative velocity are three readings of one residue, and the residue is the remainder the
linear story cannot absorb.

There are, then, two honest routes to the same reading, and it is worth naming them, because
the next two sections turn on the fact that they agree. The first route is the long way: take
the coupled measurement, take the two one-sided stories, and subtract. This is a reading
assembled out of measurements of the path, and it is the route a working instrument would
travel. The second route is the short way: read the invariant directly, as the remainder that
the construction already isolated and proved unique. The long route carries the parts and
cancels them; the short route never raises them. Both arrive at the same relative velocity,
because the whole was the parts plus the invariant all along.

That the two routes agree is the open bracket made concrete. The long route writes the reading
as a whole with its parts visibly attached; the short route writes only the remainder. The
remainder is what neither route closes --- it is exactly the quantity that survives the
subtraction --- and it is what the apparatus reads. The relative velocity is not built by the
apparatus and then measured; it is the part of the pair's motion that no story closes, and the
apparatus reads it where the stories run out. The reader should hold that picture, because the
chapter's one genuine act of selection, in the third section, is the act of certifying that the
two routes to this open quantity are routes to the \emph{same} open quantity.

\section{The Sign Is of the Interaction, Not Either Party}

The relative velocity has a magnitude and a sign, and they live in different places. This
section locates each, and the location of the sign is the section's point: the magnitude is a
property the interior energy can carry, but the sign is a property only the interaction carries.
A single party, read alone, has no relative velocity at all, and so cannot carry its sign.

Begin with what a one-sided story knows. A story tells one coordinate's motion against a fixed
background. It can be long or short, fast or slow, but it is the motion of one thing, and the
notion of relative velocity does not arise within it. Relative velocity is a reading of two
motions against each other. Hand the apparatus only the first party's story and ask for the
relative velocity, and there is nothing to read: the quantity requires a second motion to be
relative to. The relative velocity is a cross term. It is the part of the reading that couples
the two parties, and it vanishes the moment either party is removed.

This is the same blindness the book has already proved about the interior energy, now seen from
the other side. The energy the apparatus built carries magnitude and only magnitude; its
diagonal is nonnegative by construction, so a negative reading can never be a diagonal energy.
The magnitude of the relative velocity is exactly the kind of thing the energy can hold --- a
size, an amount of separation. But the sign is not a size. It is an orientation, and it lives in
the cross term: in the reading that holds both parties at once and asks which way the gap between
them is opening. The interior, holding one party's energy, is sign-blind by construction. The
cross term, holding the pair, is where the sign can be read.

So the apparatus must hold the pair to read the sign, and this is not a limitation to be
apologized for; it is the structure of what a sign of motion is. To say the pair separates in
one orientation rather than its mirror is to say something about the two parties together. It is
not a fact stored inside the first party, waiting to be retrieved, any more than it is stored
inside the second. It is a fact of the interaction, and the only instrument that can read it is
one pointed at the interaction. A reading taken of either half, however careful, returns
magnitude with the sign struck out.

The open bracket says the same thing once more. The residue sits at the center of the nest,
inside the coupling, where both parties are held together. Climb outward toward either single
party and the residue is gone: a single party closes its own story, and a closed story has no
open remainder. The relative velocity, and with it the sign, exists only at the un-closed center,
in the term that refuses to factor into one party's account and the other's. The sign is of the
interaction because the residue is of the interaction, and the residue is of the interaction
because it is precisely what neither party's story, taken alone, can close.

\section{The One Located Receipt}

The apparatus has two routes to the relative velocity, and they return the same number. That is
a finite fact, and the book does not ask the reader to take it on faith: the whole is the parts
plus the invariant, so subtracting the parts from the whole and reading the invariant directly
must agree. But there is a step between \emph{the two readings are equal} and \emph{the two
readings are one reading}, and this section is about that step, because it is the place where
the relative-velocity reading forces the construction's one genuine selection into the open.

Consider what the apparatus has after it travels both routes. It has two records. One record
carries the relative velocity as a whole with its parts attached and then cancelled; the other
carries it as the bare invariant. The two records agree on the number. But they are still two
records, two ways the same velocity was written down, and the apparatus that wants to speak of
\emph{the} relative velocity --- one reading, one class --- must do something the equality alone
does not do for it. It must declare the two records to be records of the same thing, and collapse
them into a single reading.

That declaration is a receipt. It is the apparatus issuing one certificate: \emph{these two
readings name the same relative velocity, and from here they are one}. The book has issued
receipts of a routine kind throughout --- every time it agreed with itself about which of two
records was larger, or that a difference vanished, it was certifying an identification that cost
nothing, because the records were already the same object written the same way. This receipt is
not of that kind. Here the two records are genuinely different objects --- one long, one short,
one carrying parts the other never raised --- and the apparatus chooses to treat them as one. It
is a real act, performed once, at a single located place, on a real pair of witnesses whose
agreement it has in hand.

It is worth being exact about how much this receipt does and does not claim, because the
temptation is to read it as larger than it is. It does not derive the relative velocity; the
velocity was already there, the open remainder, before any receipt. It does not certify that the
two readings are equal; that was the finite fact, settled by the arithmetic. It certifies only the
collapse: that two equal readings may be spoken of as one reading, one class, the relative
velocity of the pair. This is the single sanctioned selection the construction needs, and it is
sanctioned precisely because it is fed a real agreement, not assumed in advance. The book's method
has always been to read rather than to assume, and the receipt keeps faith with that method: it
selects only where it has been handed two witnesses that already agree.

Isolating this receipt is the chapter's reason to exist as a re-reading rather than a repetition.
The earlier chapters produced the pair and signed it; they did their work without ever needing to
name the moment at which two readings of the residue become one. The relative-velocity reading
forces that moment into the open, because it offers two visibly different routes to the same open
quantity and then asks to call the two routes one. The apparatus answers with a single receipt.
That receipt is the located act of identification on which the chapter's one class of relative
velocity rests, and the reader who has followed the open bracket will recognize it for what it is:
not a closing of the residue, but a certified agreement about how the un-closed residue is to be
read.

\section{The Convention Names Which Orientation Is Matter}

The relative velocity, now a single reading, carries a sign, and the sign splits the pair. Read in
one orientation the velocity is negative; read in the mirror orientation it is positive; and the
two, summed, cancel:
\[
  \text{one orientation} + \text{its mirror} \;=\; (-1) + (+1) \;=\; 0.
\]
This is the same cancellation the book met when it first signed the pair, and it is structural for
the same reason: a quantity read once forward and once backward subtracts from itself. The new
content of this section is not the cancellation. It is the question the cancellation makes
unavoidable: if the two orientations are mirror images that sum to nothing, what tells the
apparatus to call one of them \emph{matter}?

Nothing in the relative velocity tells it. The velocity carries magnitude and orientation; it does
not carry a label. To read the sign of the interaction is to read which way the pair separates; it
is not to read which way is the favored way. The favored way is not in the residue, was never in
the residue, and the apparatus that went looking for it in the interior found only sign-blind
magnitude. The label is a separate kind of fact, and the apparatus supplies it with a separate kind
of act.

That act is counting. The apparatus tallies orientations --- how often the reading comes back one
way against how often it comes back the other --- and asks whether one way clears the other by more
than the noise floor it has used to anchor every comparison in the book. When one orientation
dominates the tally past the floor, the apparatus takes that dominance as its empirical input and
names the dominant orientation \emph{matter}. The count is the measurement: it reads an asymmetry
in how the orientations occur. The name sits on top of the count as a convention. The reading is
the asymmetry; the label is the choice of which asymmetry to honor with the word.

It is worth keeping the two layers apart, because their relationship is the whole method of the
book in miniature. Beneath is a measurement --- a tally, a dominance past a floor, a fact the
apparatus reads and can show. Above is a convention --- the assignment of the word \emph{matter} to
the dominant orientation, an assignment the apparatus makes and names as its own rather than
pretending to derive. The count does not become a convention by being counted, and the convention
does not become a measurement by being laid on a count. The apparatus reads the asymmetry and then,
in the open, chooses the word. The deeper question --- why one orientation should dominate at
all --- it does not answer here. That question belongs to a physics beyond the finite apparatus,
and the next chapter will show that even the dominance the apparatus reads is relative to the
baseline it reads against, not absolute.

\section{Measured, Not Derived}

The chapter closes by stating its ceiling and standing under it. The split between the two
orientations is the sign of the relative-velocity invariant, measured by the apparatus. The
orientation named \emph{matter} is a convention, anchored by a tally. That is the entire claim, and
the section's discipline is to refuse every larger one that the language invites.

It invites several. A signed pair that cancels looks like an account waiting to explain why the
world is not empty; a dominance past a floor looks like a mechanism waiting to be named; the words
\emph{matter} and \emph{antimatter} arrive carrying a century of physics the finite apparatus has
not done. The apparatus declines all of it. It does not derive an asymmetry of the early universe.
It does not propose a process by which one orientation came to outnumber its mirror. It does not
claim a violation, a condition, or a mechanism. The pair, summed, cancels; the cancellation is
structural, a quantity minus itself; and the dominance the apparatus reads is an empirical input it
counts, not a result it secures. The book has been honest about exactly this kind of edge from its
first pages, and it stays honest here: it measures the sign and names the convention, and it leaves
the physics that would explain the convention to a setting it has been careful never to claim.

One image from that century is worth borrowing, not to claim it but to mark how near the shapes run.
A positron is never caught at rest wearing its name; it is caught only when it annihilates --- when it
meets its mirror and the two signed readings leave together as a pair of equal photons, flung back to
back, so that what a detector registers is not the positron but its departure. The detection is the
cancellation. The thing is confirmed only in the act of summing to nothing, and that act is the one
the apparatus has read all along: an oriented quantity differing from itself by a sign which, taken
with its mirror, comes to zero. The apparatus builds no such detector, derives nothing about the
energy those photons carry --- a number it could only copy from a table --- and claims no amplitude
for the event. It claims the structure the detector trades on, and not one thing more: that an
oriented quantity can be read, in full, in the single moment it cancels. Antimatter, read this way,
is the orientation the account knows best by watching it not matter.

So the chapter ends where its figure began, at the open bracket. The relative velocity is the
residue read at the un-closed center of the nest --- the difference the coupled measurement carries
that no story closes. The apparatus reaches that residue along two routes, certifies with a single
receipt that the two routes read one quantity, reads the sign of that quantity as the orientation in
which the pair separates, and names one orientation by a convention it anchors with a count. At no
step does it close the bracket. The relative velocity is read, not constructed; its sign is read,
not derived; its name is chosen, not found.

This is why the re-reading keeps faith with the stop rule that preceded it. The previous chapter
declined to add depth, and this chapter adds none: it does not cut deeper into the pair, does not
secure a new formal result the construction did not already hold, does not introduce a distinction
without which the finite account would collapse. It takes a reading the pair always carried and had
not been asked for --- the relative velocity, and the sign of the interaction --- and it takes that
reading by the only honest means the book has ever used. The instrument is pointed at the pair; the
open quantity is read; the sign is the orientation of the reading; and the apparatus, having
measured what it can and named what it cannot derive, leaves the bracket open, exactly as it found
it.
